\documentclass[aspectratio=169,10pt]{beamer}
\newcommand{\reporoot}{../..}
% Shared preamble of the main deck and the topic decks.
% Define \reporoot (the path from the .tex file to the repository root)
% before \input-ing this file.
\usetheme[progressbar=frametitle,numbering=fraction,sectionpage=progressbar]{metropolis}

\usepackage{amsmath,amssymb}
\usepackage{booktabs}
\usepackage{graphicx}
\usepackage{tikz}
\usetikzlibrary{arrows.meta,positioning,decorations.pathreplacing}
\usepackage[authoryear,round]{natbib}
\usepackage{xcolor}

\graphicspath{{\reporoot/figures/}}

% Every number on the slides is produced by experiments/*.py and read from here.
\input{\reporoot/results/e1_markov_book.tex}
\input{\reporoot/results/e2_best_queues.tex}
\input{\reporoot/results/e3_queue_reactive.tex}
\input{\reporoot/results/e4_hawkes.tex}
\input{\reporoot/results/e5_signature.tex}
\input{\reporoot/results/e6_propagator.tex}
\input{\reporoot/results/e7_market_making.tex}
\input{\reporoot/results/e8_ofi.tex}
\input{\reporoot/results/e9_queue_value.tex}
\input{\reporoot/results/e10_markouts.tex}
\input{\reporoot/results/e11_glft_exact.tex}
\input{\reporoot/results/e12_signal.tex}
\input{\reporoot/results/e13_lvr.tex}
\input{\reporoot/results/e14_mlofi.tex}
\input{\reporoot/results/e15_rate_estimation.tex}
\input{\reporoot/results/e16_queue_reactive_price.tex}
\input{\reporoot/results/e17_hawkes_pitfalls.tex}
\input{\reporoot/results/e18_propagator_estimation.tex}
\input{\reporoot/results/e19_latent_book.tex}
\input{\reporoot/results/e20_microprice.tex}
\input{\reporoot/results/e21_fill_estimation.tex}
\input{\reporoot/results/e22_glosten_milgrom.tex}
\input{\reporoot/results/e23_cdl_on_book.tex}

\newcommand{\E}{\mathbb{E}}
\newcommand{\Var}{\operatorname{Var}}
\newcommand{\OFI}{\mathrm{OFI}}
\newcommand{\dd}{\mathrm{d}}
\newcommand{\src}[1]{\par\vfill{\scriptsize\color{gray}#1}}
\definecolor{accent}{HTML}{0C5DA5}
\definecolor{warn}{HTML}{C0392B}
\newcommand{\hl}[1]{\textcolor{accent}{\textbf{#1}}}

\setbeamerfont{footnote}{size=\tiny}
\setbeamerfont{caption}{size=\scriptsize}

\author{Yaohui Shu}
\institute{Stony Brook University}
\date{October 2026}

% Topic decks: a references frame that lists only the works cited in the deck.
\newcommand{\deckreferences}{%
  \begin{frame}[allowframebreaks]{References}
    \tiny
    \renewcommand{\bibsection}{}
    \bibliographystyle{plainnat}
    \bibliography{\reporoot/slides/refs}
  \end{frame}}

% Topic decks: the standard "where this fits" footer on the title frame.
\newcommand{\deckseries}{Deep dive in the series \emph{Stochastic Models of the Limit Order Book}}

% Topic decks: \decknumber{n} puts the series line under the date on the title page.
\newcommand{\decknumber}[1]{\institute{Stony Brook University\hfill{\scriptsize\color{gray}Deep dive #1 of 9}}}


\title{Order Flow Imbalance and the Micro-Price}
\subtitle{What moves the price at short horizons, and the best guess of where it goes}
\decknumber{6}

\begin{document}

\maketitle

\begin{frame}{The question}
  \small
  \begin{columns}[T]
    \begin{column}{0.55\textwidth}
      Over seconds, prices move because queues at the best quotes are added to and eaten. \hl{Which summary of order flow explains the move, and which summary of the book predicts the next one?}
      \begin{itemize}
        \item \textbf{Contemporaneous}: order flow imbalance (OFI) explains price changes in the same window.
        \item \textbf{Predictive}: queue imbalance forecasts the next move; the micro-price turns it into a fair value.
      \end{itemize}
      An LP quotes around a fair value. Quoting around the mid when the micro-price is elsewhere hands edge to the next trader.
    \end{column}
    \begin{column}{0.42\textwidth}
      \begin{alertblock}{In this deck}
        \small
        \begin{enumerate}
          \item OFI: definition, linear impact, depth.
          \item Multi-level and integrated OFI.
          \item The micro-price and how to estimate it.
          \item Evidence, and the statistics of overlapping windows.
        \end{enumerate}
      \end{alertblock}
    \end{column}
  \end{columns}
\end{frame}

\begin{frame}{Order flow imbalance}
  \small
  From consecutive best-quote snapshots $n-1 \to n$, with best bid and ask prices $p_B, p_A$ and sizes $q_b, q_a$ \citep{cont2014price}:
  \[ e_n = \mathbf{1}_{\{p_{B,n}\ge p_{B,n-1}\}} q_{b,n} - \mathbf{1}_{\{p_{B,n}\le p_{B,n-1}\}} q_{b,n-1} - \mathbf{1}_{\{p_{A,n}\le p_{A,n-1}\}} q_{a,n} + \mathbf{1}_{\{p_{A,n}\ge p_{A,n-1}\}} q_{a,n-1}. \]
  New bids, cancelled asks and market buys push $e_n$ up. Summed over a window of length $\Delta t$ (s), $\OFI = \sum e_n$ (shares or orders).
  \vspace{0.3em}
  \begin{columns}[T]
    \begin{column}{0.5\textwidth}
      \textbf{Stylized model.} If every level holds $\bar q$ orders, $\OFI$ moves one quote by $\OFI/\bar q$ ticks and the mid by half as much:
      \[ \Delta S/\delta = \kappa\,\OFI + \xi, \qquad \kappa \approx \frac{1}{2\bar q}, \]
      with mid change $\Delta S$ (currency), tick $\delta$ and residual $\xi$.
    \end{column}
    \begin{column}{0.46\textwidth}
      \centering
      \includegraphics[width=0.68\textwidth]{fig7_ofi.png}

      \raggedright\footnotesize Markovian book: $\kappa = \EeightSlope$ against $1/(2\bar q) = \EeightStylized$; $R^2 = \EeightRsq$, $\EeightN$ windows of $\EeightWindow$ s.
    \end{column}
  \end{columns}
\end{frame}

\begin{frame}{Deeper levels: multi-level and integrated OFI}
  \centering
  \includegraphics[width=0.78\textwidth]{fig14_mlofi.png}
  \small
  \begin{columns}[T]
    \begin{column}{0.48\textwidth}
      \footnotesize
      \textbf{Multi-level OFI}: the same formula at each level $m$ ticks behind the best \citep{xu2019multi}. \textbf{Integrated OFI}: level OFIs divided by mean depth and combined with the first principal component weights \citep{cont2023cross}. Here the weights are $\EfourteenWeights$.
    \end{column}
    \begin{column}{0.48\textwidth}
      \footnotesize
      Out-of-sample $R^2$, $\Delta t = 10$ s ($\EfourteenLongN$ windows): best level $\EfourteenLongOne$, $\EfourteenK$ levels $\EfourteenLongAll$, integrated $\EfourteenLongInt$. At $\Delta t = 2$ s: $\EfourteenShortOne$, $\EfourteenShortAll$, $\EfourteenShortInt$. Once the best queue empties, the next level decides how far the price moves.
    \end{column}
  \end{columns}
\end{frame}

\begin{frame}{Statistics of OFI regressions}
  \small
  \begin{columns}[T]
    \begin{column}{0.5\textwidth}
      \textbf{Tautology.} OFI includes the events that move the price, so part of $R^2$ is mechanical. \citet{cont2014price} re-estimate without price-changing events: $R^2$ falls from about 65\% to the 35\% to 60\% range, still large.

      \vspace{0.3em}
      \textbf{Heteroskedasticity.} Residual variance follows depth and time of day. Use White (heteroskedasticity-consistent) errors; for regressions of estimated coefficients across windows, Newey-West errors.

      \vspace{0.3em}
      \textbf{Nonstationary coefficient.} $\kappa$ varies with depth inside the day, so estimate it on short subsamples (CKS: half an hour) or model $\kappa \propto \bar q^{-1}$.
    \end{column}
    \begin{column}{0.46\textwidth}
      \textbf{Overlapping horizons.} Forecast regressions of the mid change over $h$ on today's imbalance, with one sample per event, have strongly overlapping targets. Effective sample size is closer to $T/h$ than to the number of events. Block bootstrap or Hansen-Hodrick errors.

      \vspace{0.3em}
      \textbf{Out of sample.} Fit on the first half, test on the second; with many levels in-sample $R^2$ always rises.
    \end{column}
  \end{columns}
\end{frame}

\begin{frame}{The micro-price}
  \small
  \begin{columns}[T]
    \begin{column}{0.52\textwidth}
      State $x = (I, s)$: imbalance $I = q_b/(q_b + q_a)$ in buckets and spread $s$ (ticks). The micro-price \citep{stoikov2018micro} is the limit of the expected mid after the next $k$ mid changes:
      \[ P^{\mathrm{micro}} = M + g_1(x) + g_2(x) + \cdots, \qquad g_{i+1} = B\,g_i , \]
      with $M$ the mid (ticks). From consecutive observations:
      \[ g_1 = (I_d - Q)^{-1} R\,K, \qquad B = (I_d - Q)^{-1} T, \]
      where $Q$ holds transitions that keep the mid, $T$ those that change it, $R(x, k)$ the probability of a mid change of $k$ ticks, $K$ the vector of changes and $I_d$ the identity.
    \end{column}
    \begin{column}{0.45\textwidth}
      \textbf{Symmetrize.} Each observation also enters mirrored ($I \to 1 - I$, $\Delta M \to -\Delta M$). Then $g_i$ is odd in imbalance and the series converges.

      \vspace{0.3em}
      \textbf{Compare with the weighted mid} $I p_A + (1 - I)p_B = M + (I - \tfrac12)s$, which has no justification beyond intuition and is not a martingale.

      \vspace{0.3em}
      \textbf{Assumption.} $(I, s)$ is Markov. It is not exactly: queue sizes, not only their ratio, matter (deck 2).
    \end{column}
  \end{columns}
\end{frame}

\begin{frame}{The micro-price on a simulated book}
  \centering
  \includegraphics[width=0.86\textwidth]{fig20_microprice.png}
  \begin{columns}[T]
    \begin{column}{0.5\textwidth}
      \footnotesize
      Power-law book, $\EtwentyNevents$ events; first half trains, second half tests. (a) Model $\sum_{i\le k} g_i$ against realized mid changes after $k$ moves, block errors: in inner buckets the worst gap is $\EtwentyZmax$ s.e., a sign that $(I, s)$ is not quite Markov. $G$ reaches $\EtwentyGmax$ ticks at one-tick spread.
    \end{column}
    \begin{column}{0.46\textwidth}
      \footnotesize
      (b) Mean squared error of the future mid relative to the mid: micro-price $\EtwentyMicroOne\%$ better at 1 s, weighted mid $\EtwentyWmidOne\%$; at 10 s, $\EtwentyMicroTen\%$ and $\EtwentyWmidTen\%$. In this book $G(I)$ is nearly linear, so the two are close; real books differ (next slide).
    \end{column}
  \end{columns}
\end{frame}

\begin{frame}{What the data show}
  \small
  \begin{itemize}
    \item \textbf{OFI explains short-horizon moves.} Fifty S\&P 500 stocks drawn at random, TAQ data for April 2010, 10-second windows, half-hour regressions: average $R^2$ of 65\%. A quadratic term raises it only to 68\%. The slope is inversely proportional to depth; near the open depth is about half its average and impact about twice its average \citep{cont2014price}.
    \item \textbf{Trades alone explain less.} The relation between price changes and trade volume is noisier and less robust than the one with OFI \citep{cont2014price}.
    \item \textbf{Deeper levels help.} Multi-level OFI improves the fit over the best level \citep{xu2019multi}. Integrated OFI makes contemporaneous cross-impact between stocks redundant, while lagged cross-asset OFI forecasts short-horizon returns \citep{cont2023cross}.
    \item \textbf{Micro-price.} Estimated on high-frequency data, it predicts short-term prices better than the mid or the weighted mid \citep{stoikov2018micro}.
    \item \textbf{Deep models.} Networks on raw book states forecast short-term moves out of sample \citep{zhang2019deeplob}; deep OFI features forecast returns at several horizons \citep{kolm2023deep}; pooled training finds a universal price formation mechanism \citep{sirignano2019universal}.
  \end{itemize}
\end{frame}

\begin{frame}{Discussion and takeaways}
  \small
  \begin{columns}[T]
    \begin{column}{0.5\textwidth}
      \textbf{For a liquidity provider}
      \begin{itemize}
        \item Quote around the micro-price, not the mid. The gap is up to a few tenths of a tick at the extremes of imbalance here, comparable to half the spread.
        \item OFI is contemporaneous; imbalance is predictive. A strategy can only use the second in real time.
        \item Imbalance is a signal with a finite life. How much of it to put into quotes depends on that life (deck 7).
      \end{itemize}
    \end{column}
    \begin{column}{0.46\textwidth}
      \textbf{Open questions}
      \begin{itemize}
        \item A micro-price that uses queue sizes and recent flow, with uncertainty, and stays a martingale.
        \item OFI across fragmented venues with latency between them.
        \item Whether generative models of the book reproduce the OFI slope and its depth dependence (frontier part of the main deck).
      \end{itemize}
    \end{column}
  \end{columns}
\end{frame}

\begin{frame}[fragile]{Reproduce}
  \small
  Code: \texttt{impact.py} (OFI, multi-level, integrated), \texttt{microprice.py}. Experiments: \texttt{e8\_ofi.py}, \texttt{e14\_mlofi.py}, \texttt{e20\_microprice.py}.
\begin{verbatim}
from stochastic_lob import microprice as mp
idx, imb = mp.states(rec.q_bid, rec.q_ask, rec.best_ask - rec.best_bid)
fit = mp.fit(idx[:n // 2], mid[:n // 2])          # g_1, B, G per state
mid_, wmid, micro = mp.predictors(idx, imb, mid, spread, fit["G"])
\end{verbatim}
\end{frame}

\deckreferences
\end{document}
